\chapter{Related Work}\label{ch:related}
%% ---------------------------------------------------------------------
%% PURPOSE: position your work. Show you know the field, and make clear
%% what is new about your contribution relative to the closest work.
%%
%% Rules of thumb:
%%   - Organise by THEME or APPROACH, not paper by paper. Each section is
%%     one family of methods.
%%   - Every section ends with a contrast sentence: ``Unlike these works,
%%     we ...''. A list of summaries with no contrast is an annotated
%%     bibliography, not related work.
%%   - Be fair: state what each line of work does well before its limits.
%%   - A comparison table (approach x property) is often clearer than prose.
%%   - Cite primary sources. Check claims you attribute to others against
%%     the paper itself.
%% ---------------------------------------------------------------------

\section{\todo{Theme A}}\label{sec:related-a}
\todo{Family of approaches A. What they do; where they stop.}
Unlike these works, \todo{our difference}.

\section{\todo{Theme B}}\label{sec:related-b}
\todo{Family of approaches B.}

\section{Summary}\label{sec:related-summary}
\begin{table}[htbp]
  \centering
  \caption{Comparison of related approaches. \todo{Describe the columns.}}
  \label{tab:related-comparison}
  \begin{tabular}{lccc}
    \toprule
    Approach & Property 1 & Property 2 & Property 3 \\
    \midrule
    Work A   & \checkmark & --         & --         \\
    Work B   & --         & \checkmark & --         \\
    \textbf{This thesis} & \checkmark & \checkmark & \checkmark \\
    \bottomrule
  \end{tabular}
\end{table}
